% GENERATED FILE. Edit canonical lessons, then run npm run generate:books.
% canonical-source-hash: fnv1a64:41b758c5e8ae8c1b
% canonical-lessons: AR-C104-janub, AR-C104-sharq, AR-C104-gharb, AR-C104-makan, AR-C104-unwan

\chapter{South, East, West, Place, Address}
\label{ch:ar-south-east-west-place-address}
\hlchaptermodality{104}

\begin{chapteropening}
\textbf{By the end of this chapter:} \emph{I can say south, east, west, a place and an address.}

Complete the last lesson of chapter 104: I can say south, east, west, a place and an address.
\end{chapteropening}

\section[\texorpdfstring{\ar{جنوب} (janūb)}{جنوب (janūb)}]{\ar{جنوب} (janūb) — south}

\label{lesson:AR-C104-janub}



\begin{quote}
Before the new one: say the Arabic for a corner, then the Arabic for north.
\end{quote}

\subsection*{You'll want to know: \ar{جنوب}}
\textbf{\ar{جنوب}} — \emph{janūb} — \textquotedblleft{}south\textquotedblright{}.

\begin{cousinweb}
\textbf{\ar{جنوب}} (\emph{janūb}) is built on the native Arabic root \emph{j-n-b}, \textquotedblleft{}the side; the south wind\textquotedblright{} (Lane). Lane gives \emph{janūb} as the south wind; the same root gives \emph{jānib}, \textquotedblleft{}side\textquotedblright{}.
\end{cousinweb}

\begin{grammarlens}[title={What you've built}]
One more word for this chapter.
\end{grammarlens}

\subsection*{Your turn}
\begin{itemize}
\raggedright
  \item \emph{Say it:} \emph{janūb}
  \item \emph{Say it:} \emph{janūb}, once more
  \item \emph{Say it:} say \emph{janūb}
  \item \emph{Recall:} say the Arabic for a corner, then the Arabic for north, then say \emph{janūb} again
\end{itemize}

\begin{tcolorbox}[breakable,colback=teal!4,colframe=teal!35!black,title={Before you move on}]
What does \ar{جنوب} mean? (\textquotedblleft{}South\textquotedblright{}.) Say it once more.
\end{tcolorbox}
\section[\texorpdfstring{\ar{شرق} (sharq)}{شرق (sharq)}]{\ar{شرق} (sharq) — east}

\label{lesson:AR-C104-sharq}



\begin{quote}
Before the new one: say the Arabic for north, then the Arabic for south.
\end{quote}

\subsection*{You'll want to know: \ar{شرق}}
\textbf{\ar{شرق}} — \emph{sharq} — \textquotedblleft{}east\textquotedblright{}.

\begin{cousinweb}
\textbf{\ar{شرق}} (\emph{sharq}) is built on the native Arabic root \emph{sh-r-q}, \textquotedblleft{}to rise (of the sun)\textquotedblright{} (Lane). \emph{Sharq} is where the sun rises: the east.
\end{cousinweb}

\begin{grammarlens}[title={What you've built}]
One more word for this chapter.
\end{grammarlens}

\subsection*{Your turn}
\begin{itemize}
\raggedright
  \item \emph{Say it:} \emph{sharq}
  \item \emph{Say it:} \emph{sharq}, once more
  \item \emph{Say it:} say \emph{sharq}
  \item \emph{Recall:} say the Arabic for north, then the Arabic for south, then say \emph{sharq} again
\end{itemize}

\begin{tcolorbox}[breakable,colback=teal!4,colframe=teal!35!black,title={Before you move on}]
What does \ar{شرق} mean? (\textquotedblleft{}East\textquotedblright{}.) Say it once more.
\end{tcolorbox}
\section[\texorpdfstring{\ar{غرب} (gharb)}{غرب (gharb)}]{\ar{غرب} (gharb) — west}

\label{lesson:AR-C104-gharb}



\begin{quote}
Before the new one: say the Arabic for south, then the Arabic for east.
\end{quote}

\subsection*{You'll want to know: \ar{غرب}}
\textbf{\ar{غرب}} — \emph{gharb} — \textquotedblleft{}west\textquotedblright{}.

\begin{cousinweb}
\textbf{\ar{غرب}} (\emph{gharb}) is built on the native Arabic root \emph{gh-r-b}, \textquotedblleft{}to go away; to set (of the sun)\textquotedblright{} (Lane). \emph{Gharb} is where the sun sets; the same root gives \emph{maghrib}, the source of English \emph{Maghreb}.
\end{cousinweb}

\begin{grammarlens}[title={What you've built}]
One more word for this chapter.
\end{grammarlens}

\subsection*{Your turn}
\begin{itemize}
\raggedright
  \item \emph{Say it:} \emph{gharb}
  \item \emph{Say it:} \emph{gharb}, once more
  \item \emph{Say it:} say \emph{gharb}
  \item \emph{Recall:} say the Arabic for south, then the Arabic for east, then say \emph{gharb} again
\end{itemize}

\begin{tcolorbox}[breakable,colback=teal!4,colframe=teal!35!black,title={Before you move on}]
What does \ar{غرب} mean? (\textquotedblleft{}West\textquotedblright{}.) Say it once more.
\end{tcolorbox}
\section[\texorpdfstring{\ar{مكان} (makān)}{مكان (makān)}]{\ar{مكان} (makān) — a place}

\label{lesson:AR-C104-makan}



\begin{quote}
Before the new one: say the Arabic for east, then the Arabic for west.
\end{quote}

\subsection*{You'll want to know: \ar{مكان}}
\textbf{\ar{مكان}} — \emph{makān} — \textquotedblleft{}a place\textquotedblright{}.

\begin{cousinweb}
\textbf{\ar{مكان}} (\emph{makān}) is built on the native Arabic root \emph{k-w-n}, \textquotedblleft{}to be, to exist\textquotedblright{} (Lane). \emph{Makān} is the \emph{ma-} place noun, \textquotedblleft{}a place of being\textquotedblright{}; Lane notes some put it under \emph{m-k-n}.
\end{cousinweb}

\begin{grammarlens}[title={What you've built}]
One more word for this chapter.
\end{grammarlens}

\subsection*{Your turn}
\begin{itemize}
\raggedright
  \item \emph{Say it:} \emph{makān}
  \item \emph{Say it:} \emph{makān}, once more
  \item \emph{Say it:} say \emph{makān}
  \item \emph{Recall:} say the Arabic for east, then the Arabic for west, then say \emph{makān} again
\end{itemize}

\begin{tcolorbox}[breakable,colback=teal!4,colframe=teal!35!black,title={Before you move on}]
What does \ar{مكان} mean? (\textquotedblleft{}A place\textquotedblright{}.) Say it once more.
\end{tcolorbox}
\section[\texorpdfstring{\ar{عنوان} (ʿunwān)}{عنوان (ʿunwān)}]{\ar{عنوان} (ʿunwān) — an address}

\label{lesson:AR-C104-unwan}



\begin{quote}
Before the new one: say the Arabic for west, then the Arabic for a place.

Say \textquotedblleft{}good evening.\textquotedblright{} (\emph{Masāʾ al-khayr} — and the little \textbf{\ar{ء}} that closes \emph{masāʾ}.)
\end{quote}

\subsection*{You'll want to know: \ar{عنوان}}
\textbf{\ar{عنوان}} — \emph{ʿunwān} — \textquotedblleft{}an address\textquotedblright{}.

\begin{cousinweb}
\textbf{\ar{عنوان}} (\emph{ʿunwān}) is built on the native Arabic root \emph{ʿ-n-w}, \textquotedblleft{}the title of a letter or book\textquotedblright{} (Lane). Lane's \emph{ʿunwān} is the title or superscription of a writing; the postal address is a later sense.
\end{cousinweb}

\begin{grammarlens}[title={What you've built}]
One more word for this chapter.
\end{grammarlens}

\subsection*{Your turn}
\begin{itemize}
\raggedright
  \item \emph{Say it:} \emph{ʿunwān}
  \item \emph{Say it:} \emph{ʿunwān}, once more
  \item \emph{Say it:} say \emph{ʿunwān}
  \item \emph{Recall:} say the Arabic for west, then the Arabic for a place, then say \emph{ʿunwān} again
\end{itemize}

\begin{tcolorbox}[breakable,colback=teal!4,colframe=teal!35!black,title={Before you move on}]
What does \ar{عنوان} mean? (\textquotedblleft{}An address\textquotedblright{}.) Say it once more.
\end{tcolorbox}
